\chapter{查理听见歌声}

查理第一次听见蓝鲸歌声时，先哭了。

他甚至没有立刻明白自己为什么哭。眼前没有战场，没有尸山，没有瑟恩在控制区燃烧的翅膀，也没有渊民族老人胸前闭合的水膜。只有一片深蓝色的海，一头庞大的蓝鲸在远处缓慢下潜。

声音从海水里来。

低、长、平稳，像某种把时间压成波形的呼吸。它穿过查理的意识时，没有解释任何东西，也没有给出答案。可他的身体反应比理解更快。胸口先酸，喉咙发紧，接着眼泪就掉了下来。

在现实里，他大概还躺在灵接舱里。

在这段记录里，他没有真正的眼睛。

可他还是哭了。

因为那声音和法厄同终战中的巨鲸临终频率贴得太近。

不完全一样。蓝鲸的歌更柔软，更受身体限制，也更孤独。它少了星海档案馆里那种古老而庞大的层次，却保留了最深处的一点方向感。

像一封被海水泡了亿万年的信，字迹早已模糊，只剩开头那个称呼还在。

查理想起瑟恩最后说：

如果将来真的有什么东西从地球海里唱回来，记得替我听一下。

“我听见了。”查理对着海说。

没人回应。

蓝鲸继续下潜。它不知道自己替谁传话，也不知道自己的歌会被一个来自另一个时代、另一个意识事故里的年轻人听见。它只是游动、呼吸、歌唱。它的存在足够自然，正因为自然，才更让查理难受。

巨鲸没有复活。

法厄同也没有回来。

瑟恩、老人、那个断开的节肢族、岚柯丢掉的库、索勒烧掉的太阳环，全都没有被补偿。

可有一声歌活下来了。

查理终于明白，所谓“没有完全死”，并不等于原样归来。很多东西只能以极低、极慢、极陌生的方式继续存在。存在到后来，连它自己都不记得曾经为什么出发。

海水里，蓝鲸的长鸣再次传来。

这一次，查理听得更清楚。

那声音深处有星海的低频，有渊民族水膜的回声，有震翼族最后断路时留下的银蓝残响，还有洛安抱着容器时那种怕得要命却仍然不肯松手的心跳。

它们都很淡。

淡到任何科学仪器都只会把它们归入自然声学。

可查理知道。

知道不代表能够证明。也不代表可以让别人相信。他只是站在这段记录里，像一个迟到亿万年的听众，替那些没有抵达地球的人，把这声歌听完。

蓝鲸浮上海面，喷出一柱白色水雾。

阳光落在它背上。

查理突然很想对吉姆和安娜说点什么。他想说，我们可能不是某一族创造的，我们身上也许有太多失败者留下的痕迹。他想说，地球不是奖品，人类也不是任何文明的最终答案。

可话到最后，只剩一句很小的声音：

“瑟恩，它真的唱回来了。”

海面没有回答。

远处，蓝鲸再次沉入水下。
